\documentclass[12pt]{article}
\usepackage[utf8]{inputenc}
\usepackage[T1]{fontenc}
\usepackage[margin=1in]{geometry}
\usepackage{setspace}
\usepackage{xcolor}
\usepackage[colorlinks=true, urlcolor=blue!60!black, linkcolor=blue!60!black]{hyperref}
\setlength{\parindent}{2em}
\setlength{\parskip}{1.25em}
\renewcommand{\baselinestretch}{1.0}
\begin{document}
\begin{center}
{\LARGE \textit{Fitzpatrick v. Bitzer}, 427 U.S. 445 (1976)\par}
\vspace{0.8em}
{\large Case Brief --- Four-Page Edition\par}
\end{center}
\vspace{1em}

\section{Memory Jogger}

Because the Fourteenth Amendment expressly limits state authority and shifted the federal--state balance, Congress may use its Section 5 enforcement power to abrogate state sovereign immunity and authorize private suits for retroactive money damages and attorney's fees against unconsenting States (\href{https://supreme.justia.com/cases/federal/us/427/445/}{427 U.S. 445}).

\section{Detailed Case Facts}

Since 1939, Connecticut's State Employees' Retirement Act let female employees with 25 years of service retire on full pensions at age 50, while men with identical service had to wait until 55; sex-based rate tables also favored women. Title VII originally excluded state and local governments, but the Equal Employment Opportunity Act of 1972, expressly relying on Section 5 of the Fourteenth Amendment, made governments covered employers and authorized private suits for injunctions, retroactive backpay, and ``a reasonable attorney's fee as part of the costs.''

Garland M. Fitzpatrick and other male employees sued J. Frederick Bitzer, Chairman of the State Employees' Retirement Commission, and other officials for prospective relief, retroactive benefits, and fees. Because any award would ultimately come from general tax revenues, all parties treated retroactive benefits as a money judgment against the state treasury.

\section{Procedural History}

The District of Connecticut held that the retirement plan violated Title VII and enjoined it, but relying on the Eleventh Amendment and \textit{Edelman v.\ Jordan} (1974) denied retroactive benefits and fees. The Second Circuit affirmed the denial of backpay but reversed on fees, which it found only an ``ancillary effect'' on the treasury. The Supreme Court granted cross-petitions and decided the case on June 28, 1976.

\section{Judicial Votes}

Unanimous in the judgment, \textbf{9--0}.
\begin{itemize}
\item \textbf{Majority (7):} Justice Rehnquist, joined by Chief Justice Burger and Justices Stewart, White, Marshall, Blackmun, and Powell.
\item \textbf{Concurring in the judgment:} Justice Brennan; Justice Stevens.
\item \textbf{Dissents:} None.
\end{itemize}

\section{Holding}

The Eleventh Amendment and the state sovereignty it embodies are limited by Section 5 of the Fourteenth Amendment. When Congress enforces the Fourteenth Amendment by ``appropriate legislation'' and clearly authorizes private suits against States, it may abrogate sovereign immunity and allow retroactive damages, such as Title VII backpay, and attorney's fees from the state treasury. The Court assumed, without deciding, that Title VII's substantive rules were valid Section 5 legislation (a question later policed in \textit{City of Boerne}), and its contrast with Article I foreshadowed \textit{Seminole Tribe} (1996).

\section{Analysis of Opinions}

\textbf{Rehnquist (majority).} \textit{Edelman} barred retroactive payments because the statutes there lacked the ``threshold fact of congressional authorization''; here the 1972 Amendments clearly supplied it. The Fourteenth Amendment, ratified long after the Eleventh, couples prohibitions aimed at the States with an affirmative enforcement grant. Relying on \textit{Ex parte Virginia}, which called the Civil War Amendments ``limitations of the power of the States and enlargements of the power of Congress,'' Rehnquist concluded that Congress may authorize suits ``which are constitutionally impermissible in other contexts.'' Fees were valid under Section 5 for the same reason, without deciding whether they were merely ``ancillary.''

\textbf{Brennan (concurring in the judgment).} The Eleventh Amendment bars only suits by citizens of \emph{other} States, and the States surrendered common-law immunity ``in the plan of the Convention'' when they delegated powers to Congress, under both the Commerce Clause and Section 5.

\textbf{Stevens (concurring in the judgment).} The Commerce Clause amply supports regulating state employment, and it was unclear the Amendments were ``needed to secure the guarantees of the Fourteenth Amendment.'' Because backpay would come from ``separate and independent pension funds,'' he would have narrowed \textit{Edelman} and treated fees as ordinary costs.

\section{Examples: Future Applications}

\begin{enumerate}
\item \textbf{Same-Side (Valid): FMLA Family-Care Leave.} On a record of gender stereotypes in state leave policies, Congress lets state employees sue for damages. Sex classifications get heightened scrutiny, so abrogation is valid (as in \textit{Hibbs}, 2003).
\item \textbf{Same-Side (Valid): Civil Rights Fees.} Congress authorizes fee awards against States when plaintiffs prevail on substantial Fourteenth Amendment claims, including by settlement, directly extending \textit{Fitzpatrick}'s fee holding (as in \textit{Maher v.\ Gagne}, 1980).
\item \textbf{Opposite-Side (Invalid): Article I Abrogation.} Acting under the Indian Commerce or Commerce Clause, Congress lets tribes sue States. Article I predates the Eleventh Amendment and cannot abrogate immunity (\textit{Seminole Tribe}).
\item \textbf{Opposite-Side (Invalid): Age or Disability Damages.} Congress subjects States to damages for age or disability discrimination without a record of unconstitutional conduct. Lacking ``congruence and proportionality'' under \textit{City of Boerne}, the statute is not ``appropriate legislation.''
\item \textbf{Fence-Sitter: Disparate-Impact Liability.} A Section 5 statute lets employees sue state universities for damages over sex discrimination shown only by disparate impact. \textit{Fitzpatrick} left the question open in footnote 11, while post-\textit{Boerne} doctrine demands proportionality to actual violations.
\end{enumerate}

\section{Critique}

\textbf{Originalist and textualist views.} Many originalists praise the ``last-in-time'' logic: the 1868 Amendment restricted state sovereignty and so qualifies earlier immunity. Stricter textualists object that neither Section 1 nor Section 5 mentions the Eleventh Amendment or damages suits, and that the Court brushed past whether Title VII tracked actual constitutional violations, forcing \textit{City of Boerne}'s later test.

\textbf{Progressive and doctrinal views.} Scholars sympathetic to Brennan argue the Eleventh Amendment never barred suits by a State's own citizens, making the wall between Article I and Section 5 unsupported. After \textit{Seminole Tribe} and \textit{City of Boerne}, critics say, \textit{Fitzpatrick} became a narrow bottleneck that strips state employees of age- and disability-discrimination remedies.

\section{Key Quotations}

\begin{quote}\sloppy
``But we think that the Eleventh Amendment, and the principle of state sovereignty which it embodies, see \textit{Hans v.\ Louisiana}, 134 U.S.\ 1 (1890), are necessarily limited by the enforcement provisions of \S~5 of the Fourteenth Amendment.''
\end{quote}
\noindent --- Justice Rehnquist, Opinion of the Court, 427 U.S.\ at 456.

\begin{quote}\sloppy
``When Congress acts pursuant to \S~5, not only is it exercising legislative authority that is plenary within the terms of the constitutional grant, it is exercising that authority under one section of a constitutional Amendment whose other sections, by their own terms, embody limitations on state authority. We think that Congress may, in determining what is `appropriate legislation' for the purpose of enforcing the provisions of the Fourteenth Amendment, provide for private suits against States or state officials which are constitutionally impermissible in other contexts.''
\end{quote}
\noindent --- Justice Rehnquist, Opinion of the Court, 427 U.S.\ at 456.

\begin{quote}\sloppy
``The legislation considered in each case was grounded on the expansion of Congress' powers -- with the corresponding diminution of state sovereignty -- found to be intended by the Framers and made part of the Constitution upon the States' ratification of those Amendments, a phenomenon aptly described as a `carv[ing] out' in \textit{Ex parte Virginia}, supra, at 346.''
\end{quote}
\noindent --- Justice Rehnquist, Opinion of the Court, 427 U.S.\ at 455--56.

\end{document}
